\section{Part VIII: Customer Decisions Across Distinct Industries}

\subsection{Consumer Staples: The 3-Second Purchase Decision}

In fast-moving consumer goods, purchasing decisions occur in seconds amidst shelf clutter. Demanding elaborate rational narratives from supermarket shoppers is misleading: choice is heavily shaped by \textbf{physical availability, pack architecture, sensory habit, price tiers, and mental salience}.

Research must focus on in-store behavioral observation:
\begin{itemize}[leftmargin=1.2em,itemsep=0.2ex]
  \item Which competing brand was held in the hand and replaced on the shelf?
  \item Does the consumer buy on promotional discount or at full retail price?
  \item Most critically: does the household execute a \textbf{second and third purchase} at full price without launch incentives?
\end{itemize}
In FMCG, true customer validation is measured in unprompted repeat velocity.

\subsection{Luxury and Identity Goods: Social Signaling}

Luxury, fashion, and premium artisanal goods are not bought for utilitarian efficiency. Customers purchase \textbf{aesthetic craftsmanship, provenance, exclusivity, self-identity, and social signaling}.

An unknown brand cannot validate demand through surveys asking if an aesthetic design looks ``luxurious.'' The founder must test actual willingness to pay premium prices, patience in accepting waiting lists, gifting rituals, and how customers describe the object to peers. If the social circle does not recognize the signal, the economic premium evaporates.

\subsection{B2B Software and AI: Workflow Outcomes Over Features}

Corporate software users do not buy code repositories or algorithmic architectures; they purchase \textbf{reliable workflow transformations}.

When studying enterprise AI tools, observe the complete end-to-end task:
\begin{itemize}[leftmargin=1.2em,itemsep=0.2ex]
  \item What human labor occurs before the tool is invoked (e.g., data cleaning, prompt drafting)?
  \item What manual verification occurs after the tool generates output (e.g., cross-referencing citations, correcting hallucinations)?
  \item How does the output integrate into downstream ticketing systems or production databases?
\end{itemize}
If an AI assistant saves 5 minutes of drafting but adds 25 minutes of manual verification, it represents negative net customer value.

\subsection{Consulting and Expert Services: Judgment Under Uncertainty}

In high-end consulting, clients purchase \textbf{expert diagnostic judgment and risk reduction under deep uncertainty}. Because clients cannot evaluate work quality prior to engagement, decisions are dominated by reputation, intellectual rigor, peer references, and method transparency.

Discover the acute trigger that made internal inaction impossible: a boardroom audit, a regulatory subpoena, or an unexpected technological disruption. An emerging firm breaks into this market not through generic pitch decks, but by offering a \textbf{tightly bounded, highly rigorous diagnostic} that demonstrates undeniable expertise on a small, non-threatening scope before proposing large transformations.

\subsection{Two-Sided Marketplaces: The Liquidity Trilemma}

Marketplaces must understand 2 distinct customer groups simultaneously: supply and demand. A buyer seeks liquidity, trust, pricing transparency, and recourse. A seller seeks steady transaction volume, prompt payment, and platform fairness.

A marketplace can fail in 3 distinct ways:
\begin{enumerate}[leftmargin=1.2em,itemsep=0.2ex]
  \item \textbf{Search failure:} Buyers search but find inadequate local density.
  \item \textbf{Quality failure:} The transaction takes place, but service quality disappoints.
  \item \textbf{Disintermediation failure:} Buyer and seller meet on the platform, then bypass the platform for all subsequent repeat transactions.
\end{enumerate}
Discovery must map the entire transaction loop within a single \textbf{atomic market} (a tight geographic catchment or narrow professional niche) to prove repeatable matching density before expanding.

\subsection{Credence Services: Healthcare, Finance, and Legal}

In credence goods, customers cannot accurately evaluate quality even after consuming the service (e.g., medical diagnostics, financial asset management, or corporate legal structuring). Choices are governed by institutional credentials, state licenses, fiduciary duties, and institutional reputations.

Entrepreneurs in these spaces must rigorously map the tripartite division between \textbf{the patient/user, the provider/clinician, and the payer/insurer}. Never infer commercial feasibility from consumer desire alone without verifying regulatory authorization and reimbursement pathways.

\subsection{Local Services and Industrial Supplies: Uptime and Reliability}

In local trades, specialized equipment repair, and industrial input supplies, competitive advantage rarely stems from novel conceptual features. It stems from \textbf{operational uptime, geographic proximity, and delivery reliability}.

An industrial component supplier wins contracts because their delivery truck arrives within 90 minutes of an assembly line breakdown, accompanied by verified metallurgical testing certificates. The purchasing manager is buying operational peace of mind and production continuity.

\subsection{Education as a Transformation and Trust Market}

Education is a unique economic good: it is simultaneously an investment in future capability, a credential signaling device, and an intimate personal transformation. Crucially, \textbf{education is co-produced}: the institution can provide world-class instruction, rigorous feedback, and inspiring problems, but if the student does not invest cognitive struggle and disciplined hours, capability will not emerge.

A founder must uncover:
\begin{itemize}[leftmargin=1.2em,itemsep=0.2ex]
  \item What baseline prerequisites must the learner possess for instruction to take hold?
  \item What structural mechanisms keep the learner engaged when motivation collapses in week 4?
  \item What credible evidence will convince external employers that the learner's capability is authentic?
\end{itemize}
Selling an educational promise without verifying whether students can sustain the requisite cognitive labor is an ethical and commercial disaster.

\subsection{Technological Migration: When Content Commoditizes}

When generative AI and open internet repositories make explanations, basic code generation, and lecture summaries virtually free and universally accessible, customer demand migrates rapidly.

Customers no longer pay for access to recorded information. They pay for:
\begin{itemize}[leftmargin=1.2em,itemsep=0.2ex]
  \item \textbf{Rigorous, unsparing assessment:} Independent verification of whether code is production-ready.
  \item \textbf{Authentic, complex problem spaces:} Problems with noisy data and conflicting trade-offs that cannot be solved by naive prompts.
  \item \textbf{Accountability and human review:} Experienced engineering critique that uncovers deep architectural flaws.
\end{itemize}
As underlying technologies evolve, a founder must continuously reassess what remains scarce and valuable in the customer's learning journey.
